\documentclass[11pt,a4paper]{article}
\usepackage[utf8]{inputenc}
\usepackage[T1]{fontenc}
\usepackage[french]{babel}
\usepackage{lmodern}
\usepackage{geometry}
\usepackage{amsmath,amssymb}
\usepackage{booktabs}
\usepackage{enumitem}
\usepackage{hyperref}
\geometry{margin=25mm}
\hypersetup{colorlinks=true,linkcolor=blue,citecolor=blue,urlcolor=blue}
\setlength{\parskip}{0.55em}
\setlength{\parindent}{0pt}

\title{Pistes bibliographiques pour estimer l'effet de ressort}
\author{Note de travail -- méthodes alternatives et complémentaires}
\date{17 novembre 2026}

\begin{document}
\maketitle

\begin{abstract}
Cette note rassemble des approches repérées dans la littérature locale pour
estimer ou caractériser l'elastic follow-up (EFU), sans modifier le manuscrit
principal. Trois voies sont examinées : un modèle analytique à deux barres
élastoplastiques, une mesure sécante extraite d'un calcul éléments finis (EF),
et une définition multiaxiale fondée sur les taux de travail. Ces facteurs ne
sont pas interchangeables avec \(r_M\) ou \(g\) du RCC-MRx. La note précise les
hypothèses, le domaine d'emploi et une démarche de comparaison possible pour le
coude étudié.
\end{abstract}

\section{Modèle à deux barres avec plasticité dans les deux barres}

Wang, Smith et Truman étudient deux barres en série, de longueurs \(L_i\) et
sections \(A_i\), soumises à une compatibilité de déplacement et à l'équilibre.
Pour chacune, la déformation totale est représentée par
\[
  \varepsilon_i=\frac{\sigma_i}{E_i}+f_i(\sigma_i,T_i).
\]
Le facteur de suivi plastique \(Z_P\) est défini dans leur modèle à partir des
états initial et final. Lorsque les deux barres plastifient et que la loi
plastique est une loi puissance de type Ramberg--Osgood,
\[
  \varepsilon_i=\frac{\sigma_i}{E_i}+D_i\sigma_i^{m_i},
\]
leur équation (24) donne
\[
  Z_P=
  \frac{1+1/\alpha}
       {1+\dfrac{L_2}{L_1}\dfrac{D_2}{D_1}
       \left(\dfrac{A_1}{A_2}\right)^{m_2}
       \sigma_1^{\,m_2-m_1}},
  \qquad
  \alpha=\frac{K_2}{K_1}
  =\frac{A_2E_2L_1}{A_1E_1L_2}.
\]
Les paramètres \(D_i,m_i\) peuvent différer entre matériaux ou températures.
Si les deux barres sont de même matériau et à la même température, l'équation
(25) se réduit à
\[
  Z_P=
  \frac{1+(L_2/L_1)(A_1/A_2)}
       {1+(L_2/L_1)(A_1/A_2)^m}.
\]
Cette relation montre explicitement que le facteur dépend de la géométrie et
de la réponse constitutive; quand les deux barres sont plastifiées, il peut
aussi dépendre du niveau de chargement. Le cas différent où seule la barre 1
plastifie et la barre 2 reste élastique donne, dans le même article,
\(Z_P=1+1/\alpha\), indépendant de la loi plastique.

\textbf{Intérêt pour l'étude en cours.} Cette formulation permettrait de
construire un modèle réduit combinant une zone locale et la compliance du reste
de la tuyauterie, puis d'étudier l'influence de l'exposant de la loi
constitutive. Elle constitue une piste analytique distincte de la formule
RCC-MRx.

\textbf{Limites.} Le modèle est uniaxial et sériel : il ne représente ni la
flexion multiaxiale d'un coude, ni la redistribution spatiale des contraintes.
L'exposant \(m\) de la forme \(\varepsilon_p=D\sigma^m\) doit être relié avec
précaution à celui de la loi utilisée dans le calcul actuel. Faire varier
l'exposant tout en maintenant la limite conventionnelle \(R_{p0.2}\) impose
également de recalibrer \(D\); sinon on change simultanément la courbe
matériau et le facteur étudié. Une limite parfaitement plastique n'est donc
pas obtenue par simple substitution dans cette formule.

\section{Facteur sécant mesuré directement par éléments finis}

Le même article définit un facteur plastique à partir de la déformation
finale, de la contrainte finale et de la réponse élastique de référence pour le
même chargement ou déplacement. Dans sa notation uniaxiale,
\[
  Z_P =
  \frac{\varepsilon_{P,\mathrm{fl}}-\varepsilon_{P,\mathrm{eq,el}}}
       {\varepsilon_{P,\mathrm{in}}-\varepsilon_{P,\mathrm{eq,el}}},
  \qquad
  \varepsilon_{P,\mathrm{eq,el}}=\frac{\sigma_{P,\mathrm{fl}}}{E}.
\]
Cette définition fournit une mesure sécante a posteriori : elle caractérise
l'effet de ressort dans un état donné, mais ne prédit pas cet état sans calcul
inélastique.

Pour le coude, une adaptation pourrait être évaluée à chaque point d'intégration
sur la section critique, en comparant le calcul non linéaire et un calcul
élastique de référence soumis au même déplacement imposé. Il faudra définir
explicitement la grandeur scalaire comparée (composante axiale, courbure ou
mesure équivalente), l'instant de référence, et le traitement des changements
de signe. La formule publiée est établie pour le modèle uniaxial à deux barres;
son application multiaxiale ne doit pas être présentée comme une définition
universelle sans validation.

\textbf{Usage recommandé.} En faire une mesure locale de diagnostic dans le
modèle coque, et non une substitution directe à \(r_M\) ou au coefficient
d'abattement \(g\). Elle pourrait vérifier si le modèle réduit capture les
tendances observées dans le coude.

\section{Extension multiaxiale fondée sur les taux de travail}

Messner, Jetter et Sham proposent une définition locale calculable à partir
d'une simulation EF évolutive. Leur article définit, dans le cas uniaxial,
\[
  q=-\frac{\dot W_{\mathrm{inel}}}{\dot W_{\mathrm{el}}},
\]
et l'étend aux états multiaxiaux en utilisant les puissances de travail
\(\dot W_{\mathrm{inel}}=\boldsymbol{\sigma}:\dot{\boldsymbol{\varepsilon}}_{\mathrm{inel}}\)
et
\(\dot U=\boldsymbol{\sigma}:\dot{\boldsymbol{\varepsilon}}\). Lorsque la
dissipation du modèle provient uniquement de l'inélasticité, leur équation
(13) s'écrit
\[
  q=-\frac{\dot W_{\mathrm{inel}}}{\dot U-\dot W_{\mathrm{inel}}}.
\]
Cette mesure tangentielle dépend du temps et du trajet de relaxation. Elle
permet de traiter des géométries complexes et des lois constitutives réalistes,
mais nécessite un calcul EF fournissant les champs de taux de déformation et de
travail.

Les auteurs indiquent que la plasticité doit être incluse dans l'inélasticité
pour leurs applications, notamment lorsque creep et plasticité sont couplés à
haute température. Toutefois, l'article est marqué \emph{DRAFT}. La définition
est discutée principalement pour des situations de relaxation/fluage; pour un
chargement monotone quasi statique sans maintien, un facteur tangent dépend
fortement de l'incrément et ne se compare pas directement au facteur sécant
précédent. Il faudrait donc préciser le trajet, l'intervalle temporel ou
l'incrément retenu et vérifier la stabilité numérique avant toute utilisation
quantitative.

\section{Autres résultats, à ne pas confondre avec une formule générale}

\begin{itemize}[leftmargin=1.5em]
\item Takakura et al. rapportent des facteurs \(q\) obtenus par analyses EF
  pour des composants représentatifs du guide japonais à haute température.
  Les résultats dépendent de la géométrie, du chargement, du modèle
  élastoplastique et du cycle; ils peuvent servir de points de comparaison ou
  de bornes pour les cas étudiés, pas de facteur transférable sans justification
  au coude présent.
\item Messner et al. rappellent que des valeurs de conception de l'ordre de
  \(2\)--\(3\) sont utilisées dans certains cadres, mais que des concentrations
  locales peuvent conduire à des valeurs plus élevées. Il s'agit d'une
  indication de conception, non d'une prédiction pour notre géométrie.
\item Les actes UK de 1982 présentent des facteurs d'amplification de
  déformation issus de courbes fonction de la limite d'élasticité, du rapport
  des modules plastique et élastique et de la biaxialité. Ces courbes visent
  l'estimation des déformations inélastiques sous choc thermique à partir d'un
  calcul élastique; elles constituent une approche voisine, mais pas une
  définition générale de l'EFU.
\end{itemize}

\section{Proposition d'application au coude}

La voie la plus informative serait de comparer deux mesures ayant des rôles
distincts :
\begin{enumerate}[leftmargin=1.7em]
\item construire un modèle réduit zone critique--reste élastique, puis calculer
  \(Z_P\) avec la formule à deux barres pour une loi puissance, en calibrant
  chaque loi afin de conserver le même \(R_{p0.2}\);
\item extraire du modèle coque un facteur sécant local fondé sur la définition
  de Wang et al., avec la même convention de déformation et le même déplacement
  de référence;
\item si une phase de maintien/relaxation est étudiée, calculer séparément
  \(q(t)\) par le bilan de travail multiaxial de Messner et al., en vérifiant
  sa sensibilité au maillage et à l'incrément.
\end{enumerate}

Les résultats devront rester identifiés par leurs symboles et définitions :
\(Z_P\), \(q\), \(r_M\) et \(g\) ne sont pas synonymes. Une comparaison
numérique entre eux n'a de sens qu'après avoir aligné chargement, quantité
mesurée, état de référence et hypothèses constitutives.

\section*{Références}
\begin{enumerate}[leftmargin=1.7em]
\item Y. Wang, D. J. Smith et C. E. Truman, ``Inelastic Deformation and
  Elastic Follow-Up,'' \emph{Proceedings of the ASME 2013 Pressure Vessels and
  Piping Conference}, PVP2013-97744, Paris, 2013.
  \href{https://doi.org/10.1115/PVP2013-97744}{doi:10.1115/PVP2013-97744}.
\item M. C. Messner, R. I. Jetter et T.-L. Sham, ``A Method for Directly
  Assessing Elastic Follow Up in 3D Finite Element Calculations,''
  \emph{Proceedings of the ASME 2019 Pressure Vessels and Piping Conference},
  PVP2019-93644, 2019. Copie locale marquée \emph{DRAFT}.\par
  \href{https://doi.org/10.1115/PVP2019-93644}{doi:10.1115/PVP2019-93644}.
\item K. Takakura et al., ``Development of elevated temperature structural
  design guide for Japanese demonstration breeder reactor,''
  \emph{Nuclear Engineering and Design}, vol. 155, 1995, pp. 559--570.
  \href{https://doi.org/10.1016/0029-5493(94)00926-0}
  {doi:10.1016/0029-5493(94)00926-0}.
\item I. W. Goodall, D. A. Jobson et R. T. Rose, ``UK Contribution to the
  IWGFR Specialists Meeting on Advances in Structural Analysis for LMFBR
  Applications,'' Paris, 11--15 octobre 1982, document local XA0201718.
\end{enumerate}

\end{document}
